\textbf{8.\ \textsc{Brightness of Planets} (9 points)} --- \textit{Topi Löytäinen, Jaan Kalda.}

In this problem, the time periods are measured in years (y), and the distances
are measured in astronomical units (au); the distance of Earth from Sun is
$R_\oplus = 1\,\mathrm{au}$. The planets and Moon are assumed to move in the
same plane, the ecliptic, around circular orbits in the same direction.
Distance of Venus from Sun is \(R_{\text{\fontspec{DejaVuSans.ttf}♀}} = 0.72\,\mathrm{au}\), and distance
of Mars from Sun is \(R_{\text{\fontspec{DejaVuSans.ttf}♂}} = 1.5\,\mathrm{au}\). In some parts below, you
will need to assume that the planets are so-called Lambertian scatterers, i.e.\
when you can see the full disc of a planet illuminated by sunlight, middle parts
of the disc seem exactly as bright as the parts at the edge of the disc.
Throughout the problem, assume that the observer is situated in Tallinn, at the
geographical coordinates $59.5^\circ$N and $24.7^\circ$E. The axis of Earth is
tilted by $23.5^\circ$with respect to the normal of the ecliptic plane.

\textbf{i)} \textit{(0.7 points)} When can you see a waxing crescent moon?
Select one or more options (A, B and/or C), and motivate your answer with a
diagram depicting Sun, Moon and Earth. Select from the options\\
A: immediately after sunset;\\
B: during midnight;\\
C: immediately before sunrise.

\textbf{ii)} \textit{(1.2 points)}What is the culmination angle (i.e.\ the
maximal angle above the horizon) of a full moon during winter solstice in
Tallinn?

\textbf{iii)} \textit{(1.2 points)}The apparent brightness (illuminance at
Earth, i.e.\ luminous flux density) of planets may vary over a wide range of
values. By how many times does the apparent brightness of Mars vary between
when it's closest and farthest away from Earth?

\textbf{iv)} \textit{(1.2 points)} What is the time period on Earth between
the moments when Mars is the closest to Earth and farthest?

\textbf{v)} \textit{(1.2 points)}From autumn to spring, Venus can never be
seen at midnight. For how long maximally it can be seen after the sunset?

\textbf{vi)} \textit{(2.5 points)} Express the normalised apparent brightness
$I/I_0$ as a function of $R_\oplus$, \(R_{\text{\fontspec{DejaVuSans.ttf}♀}}\), and the distance $L$ from
Earth to Venus. Here, the normalisation constant $I_0$ can be chosen
arbitrarily. \textit{Hint:} $I/I_0$ should be a polynomial of $L^{-1}$.

\textbf{vii)} \textit{(1 point)} Find the distance $L = L_0$ when the apparent
brightness of Venus is the highest, and the angular distance between Venus and
Sun at that moment.
